\lysh{34843}{2}

\begin{alist}
    \item Σύμφωνα με το διάγραμμα συχνοτήτων συμπληρώνουμε τον παρακάτω πίνακα σχετικών συχνοτήτων, γνωρίζοντας ότι $f_{i} = \frac{v_{i}}{\nu}$ και $f_{i}\% = 100 f_{i}$. Είναι
    \[ \nu = \nu_1 + \nu_2 + \nu_3 = 10 + 25 + 15 = 50. \]
    \begin{center}
        \begin{mytblr}{}
            
            Βαθμολογία & Συχνότητα & Σχετική συχνότητα & Σχετική συχνότητα \\
            $x_{i}$ & $v_{i}$ & $f_{i}$ & $f_{i}\%$ \\
            
            12 & 10 & 0,2 & 20 \\
            
            13 & 25 & 0,5 & 50 \\
            
            14 & 15 & 0,3 & 30 \\
            
            Σύνολο & 50 & 1 & 100 \\
            
        \end{mytblr}
    \end{center}
    
    \item Έχουμε
    \[ \bar{x} = \frac{\sum_{i=1}^{\kappa} x_{i} \cdot v_{i}}{\nu} = \frac{12 \cdot 10 + 13 \cdot 25 + 14 \cdot 15}{50} = \frac{655}{50} = 13,1 \]
\end{alist}
\vspace{6mm}
\noindent\rule{\textwidth}{0.4pt}
\vspace{6mm}

